\section{Background and Significance}
\label{sec:background}

% Para 1: clinical motivation
Cancer knowledge bases such as the Clinical Interpretation of
Variants in Cancer (CIViC) \citep{griffith2017civic} depend on
expert curation to keep pace with growing oncology literature.
Like OncoKB \citep{chakravarty2017oncokb} and CancerMine
\citep{lever2019cancermine}, CIViC records curator-validated
drug, gene, variant, and disease assertions with their supporting
evidence; manual curation, however, lags the rate at which new
publications are indexed in PubMed \citep{civicfact2025}. Pretrained
biomedical encoders, most prominently PubMedBERT
\citep{gu2021pubmedbert} and BioLinkBERT
\citep{yasunaga2022biolinkbert}, are increasingly used to flag
candidate assertions for curator review at scale.

% Para 2: the unexamined assumption + evaluation literature gap
Models are typically selected by their performance on public
relation-extraction benchmarks: macro-averaged F1 score (macro-F1)
on BioRED \citep{luo2022biored}, BC5CDR drug--disease F1
\citep{li2016bc5cdr}, and component scores in suites such as the
Biomedical Language Understanding and Reasoning Benchmark
(BLURB) \citep{gu2021pubmedbert}. The implicit assumption is that
benchmark performance predicts downstream KB yield. To our
knowledge, this assumption has not been tested on a
curator-validated oncology knowledge base. Construct-validity
audits \citep{bowman2021construct} and evidence-centred benchmark
designs \citep{liu2024ecbd} have shown that benchmark improvements
need not imply meaningful capability gains, but those critiques
address general-domain language understanding rather than
biomedical curation.

% Para 3: biomedical NLP gap
Within biomedical NLP, benchmark suites (BLURB
\citep{gu2021pubmedbert}; BigBIO \citep{fries2022bigbio}) and
KB-curation pipelines (CIViCmine \citep{lever2019civicmine} for
variant mining; BioREx \citep{lai2023biorex} for multi-corpus
relation extraction) have developed in parallel. No study in
either body of work has tested whether benchmark-selected pipelines
optimise downstream KB yield in oncology.

% Para 4: LLM-era context and encoder-only scoping
Decoder-only large language models are now competitive with
fine-tuned encoders on biomedical relation tasks
\citep{chen2025nlpbench} and have been evaluated on BioRED and
CIViC-Fact \citep{civicfact2025}. We restricted the study to
encoder-only models because their inference pipeline and output
distribution are homogeneous across the four encoder families
compared here, so a variance decomposition can isolate encoder,
schedule, and schema effects without confounding from prompt
design or decoding configuration. The architectural diversity of
current LLMs is a separate design axis that the present factorial
cannot identify; we mark its extension as a stated follow-up
(\S\ref{sec:discussion}).

% Para 5: contributions and RQs
To close this gap, we ran a pre-registered factorial study of 310
training runs across four encoders, three nested relation schemas,
and three training schedules, evaluated against 162
curator-validated CIViC assertions. Benchmark and KB performance
are treated as two evaluation axes to be compared rather than
merged onto a single ranking
\citep{lawlor2017triangulation,liugaunt2024asq}, and the
contribution of each design lever is partitioned across the two
axes by a variance decomposition. Four research questions follow:
whether a finer relation schema improves KB surfacing (RQ1);
whether encoder scale, multi-corpus training, or
oncology-projected staging improves generalisation (RQ2); whether
encoder rankings are stable across three KB-audit metrics (RQ3);
and whether benchmark and KB rankings diverge under matched
configurations (RQ4).
